\documentclass[11pt,a4paper]{amsart}
\usepackage[T1]{fontenc}
\usepackage{lmodern}
\usepackage{amssymb}
\usepackage{microtype}
\usepackage[hyphens]{url}
\usepackage[colorlinks=true,linkcolor=black,citecolor=black,urlcolor=black]{hyperref}
\usepackage[margin=2.5cm]{geometry}

\hypersetup{
  pdftitle={Colouring the p-adic plane},
  pdfauthor={Sergi Gal\'an},
  pdfsubject={The Hadwiger-Nelson problem over the p-adic numbers},
  pdfkeywords={Hadwiger-Nelson problem, p-adic numbers, unit-distance graph, Borel chromatic number,
  measurable chromatic number, Kloosterman sums}}

\newtheorem{theorem}{Theorem}
\newtheorem{corollary}[theorem]{Corollary}
\newtheorem{proposition}[theorem]{Proposition}
\newtheorem{lemma}[theorem]{Lemma}
\theoremstyle{remark}
\newtheorem*{remark}{Remark}
\newtheorem*{question}{Question}

\newcommand{\F}{\mathbb{F}}
\newcommand{\Q}{\mathbb{Q}}
\newcommand{\R}{\mathbb{R}}
\newcommand{\Z}{\mathbb{Z}}
\newcommand{\C}{\mathbb{C}}
\newcommand{\Oo}{\mathcal{O}}
\newcommand{\Tr}{\operatorname{Tr}}
\newcommand{\N}{\operatorname{N}}
\newcommand{\chiB}{\chi_{\mathrm{B}}}
\newcommand{\chiM}{\chi_{\mathrm{m}}}

\title{Colouring the $p$-adic plane}
\author{Sergi Gal\'an}
\thanks{These results were obtained with the assistance of an AI system (Claude, by Anthropic), which carried
out the computations and the checks, and helped to draft this note. No mathematician has checked the note yet.
The author is not a professional mathematician and submits it for independent verification.}
\date{Draft, 2 October 2026}
\subjclass[2020]{Primary 05C15; Secondary 11S80, 11L05, 05C63, 52C10}
\keywords{Hadwiger--Nelson problem, $p$-adic numbers, unit-distance graph, Borel chromatic number, measurable
chromatic number, Kloosterman sums}

\begin{document}

\begin{abstract}
Let $\chi(\Q_p^2)$ be the chromatic number of the graph on $\Q_p^2$ in which $(x,y)$ and $(x',y')$ are adjacent
when $(x-x')^2+(y-y')^2=1$. Bardestani and Mallahi-Karai showed that its Borel chromatic number is finite
exactly when $p=2$ or $p\equiv 3\pmod 4$, and asked whether it is bounded independently of $p$. It is not: for
$p\equiv3\pmod4$ every proper colouring with Haar-measurable colour classes uses at least
$1+(p+1)/(2\sqrt p)$ colours. The proof is the Delsarte--Hoffman bound on the compact group $\Z_p^2$, where the
Fourier coefficients of the unit circle are Kloosterman sums at the first level and at most $2/(p+1)$ at every
deeper level. We also show $\chi(\Q_7^2)=4$, for ordinary and for Borel colourings, while $\chi(\Q_2^2)=2$ and
$\chi(\Q_3^2)=3$ follow from results of Madore. A theorem of Davies shows that the graph of every isotropic form
over a field of characteristic $0$ has infinite chromatic number, so the dichotomy of Bardestani and
Mallahi-Karai holds for ordinary colourings too; in particular $\chi(\Q_p^2)=\infty$ for $p\equiv1\pmod4$. Finally,
unit-distance graphs over $27$ real quadratic fields give $\chi(\Q_p^2)\ge4$ for every prime $p\equiv3\pmod4$ with
$7\le p<2\,129\,503\,819$, and a Chebotarev argument shows that no finite set of number fields gives this bound
for every such $p$.
\end{abstract}

\maketitle

\section{Introduction}

For a field $F$ let $G(F)$ be the graph on $F^2$ in which $v$ and $w$ are adjacent when $Q(v-w)=1$, with
$Q(x,y)=x^2+y^2$, and write $\chi(F^2)$ for its chromatic number. For $F=\R$ this is the Hadwiger--Nelson
problem, and $5\le\chi(\R^2)\le7$: the lower bound is due to de Grey \cite{deGrey} and to Exoo and Ismailescu
\cite{EI}, the upper bound to Isbell (see \cite{Soifer}). Woodall showed $\chi(\Q^2)=2$ \cite{Woodall}, and the
planes over number fields have been studied by Johnson \cite{Johnson1987}, Fischer \cite{Fischer1990}, Madore
\cite{Madore}, Moorhouse \cite{Moorhouse} and others. Madore observed that reduction modulo a prime gives
$\chi(\Q_p^2)\le\chi(\F_p^2)$ for $p\equiv3\pmod4$ \cite[\P1.5 and Remark~3.6]{Madore}.

Bardestani and Mallahi-Karai \cite{BMK} studied the graphs of quadratic forms over local fields. For a local
field $F$ of characteristic $0$ they proved that the \emph{Borel} chromatic number of the graph of a
non-degenerate form is infinite if and only if the form is isotropic, using the spectral bound of
Bachoc, DeCorte, Oliveira and Vallentin \cite{BDOV}. Over $\Q_p$ the form $x^2+y^2$ is anisotropic exactly when
$-1$ is not a square, that is for $p=2$ and $p\equiv3\pmod4$. For $p\equiv3\pmod4$ they noted that the Borel
chromatic number $\chiB(\Q_p^2)$ is $O(p^2)$, and asked:

\begin{question}[{\cite[Question~1]{BMK}}, the $p$-adic Hadwiger--Nelson problem]
Is there a constant $B$ with $\chiB(\Q_p^2)\le B$ for every prime $p\equiv3\pmod4$?
\end{question}

Write $\chiM$ for the measurable chromatic number (colour classes measurable for the Haar measure), so that
$\chi\le\chiM\le\chiB$. Our first result answers the question in the negative. We have not found another answer
in the literature, including the later papers of the same authors \cite{BMK2,BMK3}.

\begin{theorem}\label{thm:A}
For every prime $p\equiv3\pmod4$,
\[
  1+\frac{p+1}{2\sqrt p}\;\le\;\chiM(\Q_p^2)\;\le\;\chiB(\Q_p^2)\;\le\;\chi(\F_p^2),
\]
and $\chi(\F_p^2)\le(p+1)/2$ for $p>3$. In particular $\chiB(\Q_p^2)\to\infty$.
\end{theorem}

Here $\F_p^2$ denotes the graph $G(\F_p)$. Le Anh Vinh \cite[Theorem~1]{Vinh} showed
$(\tfrac12+o(1))\sqrt q\le\chi(\F_q^2)\le(q+q/r)/2$ for powers $q>3$ of an odd prime $r$; the upper bound in
Theorem~\ref{thm:A} is his. The lower bound of Theorem~\ref{thm:A} has the same leading term as his, but it holds
for all measurable colourings of the $p$-adic plane, not only for those that factor through $\F_p^2$. The exact
Fourier coefficients give better values for small $p$ (Table~\ref{tab:meas}): at least $5$ colours from $p=23$,
$6$ for $p=59$ and every $p\ge67$, $7$ for $p=71$ and every $p\ge103$.

For the ordinary chromatic number:

\begin{theorem}\label{thm:B}
$\chi(\Q_7^2)=\chiB(\Q_7^2)=4$. Moreover $\chi(\Q_2^2)=\chiB(\Q_2^2)=2$ and $\chi(\Q_3^2)=\chiB(\Q_3^2)=3$.
\end{theorem}

The upper bounds are Madore's reductions \cite[Propositions~3.2 and 3.8]{Madore}, which give locally constant
colourings, and the values $2$ and $3$ follow at once from his results (his $9$-cycle over $\Q(\sqrt7)$
\cite[Proposition~4.3]{Madore} lies in $\Q_3^2$). The new value is $\chi(\Q_7^2)=4$. Since $3$ is not a square
modulo $7$, the plane $\Q_7^2$ contains no unit equilateral triangle (the third vertex of a unit triangle on
$0$ and $(a,b)$ is $(\frac a2-\frac{\sqrt3}2b,\frac b2+\frac{\sqrt3}2a)$, which forces $\sqrt3$ into the field),
so the classical $4$-chromatic graphs, built from triangles, do not embed; we use a triangle-free $76$-vertex graph
over $\Q(\sqrt{11})$ with no proper $3$-colouring \cite{QP}, as $11\equiv2^2\pmod7$. The three ordinary values are
also proved in Lean~4 with Mathlib \cite{Lean,mathlib}, in \texttt{lean/PadicPlanes.lean} of \cite{Repo}; for
instance \texttt{(QuadraticPlanes.sumSqGraph $\Q$\_[7]).chromaticNumber = 4}.

In the isotropic case the ordinary chromatic number is infinite as well, by a theorem of Davies \cite{Davies}:

\begin{proposition}\label{prop:C}
Let $F$ be a field of characteristic $0$ and $Q$ a non-degenerate isotropic quadratic form on $F^n$, $n\ge2$.
Then the graph on $F^n$ in which $v,w$ are adjacent when $Q(v-w)=1$ has infinite chromatic number. In particular
$\chi(F^2)=\infty$ when $-1$ is a square in $F$, and $\chi(\Q_p^2)=\infty$ for $p\equiv1\pmod4$.
\end{proposition}

\begin{corollary}\label{cor:dich}
Let $F$ be a local field of characteristic $0$ and $Q$ a non-degenerate quadratic form in at least two
variables over $F$. Then the chromatic number of its graph is finite if and only if $Q$ is anisotropic, if and
only if its Borel chromatic number is finite. In particular $\chi(\Q_p^2)$ is finite exactly when $p=2$ or
$p\equiv3\pmod4$.
\end{corollary}

\begin{proof}
If $Q$ is anisotropic, \cite[Theorem~1.1]{BMK} gives $\chi\le\chiB<\infty$; if it is isotropic,
Proposition~\ref{prop:C} gives $\chiB\ge\chi=\infty$.
\end{proof}

Finally, every number field $K$ with an embedding $K\hookrightarrow\Q_p$ gives $\chi(\Q_p^2)\ge\chi(K^2)$.
Unit-distance graphs with no proper $3$-colouring over $27$ real quadratic fields \cite{QP} give:

\begin{theorem}\label{thm:D}
$\chi(\Q_p^2)\ge4$ for every prime $p\equiv3\pmod4$ with $7\le p<2\,129\,503\,819$. Moreover
$\chi(\Q_{83}^2)\ge5$.
\end{theorem}

No finite set of number fields can do this for every $p$:

\begin{proposition}\label{prop:E}
Let $\mathcal K$ be a finite set of number fields with $\chi(K^2)\ge4$ for every $K\in\mathcal K$. Then there are
infinitely many primes $p\equiv3\pmod4$ such that no $K\in\mathcal K$ embeds in $\Q_p$.
\end{proposition}

For the $27$ fields of Theorem~\ref{thm:D} the first such prime after $3$ is $2\,129\,503\,819$. The proposition
is about graphs over fixed fields; a single abstract graph might be realised in $\Q_p^2$ over different fields for
different $p$.

\section{Reduction modulo \texorpdfstring{$p$}{p}: upper bounds}

\begin{lemma}[Madore {\cite[Propositions~3.2 and 3.8]{Madore}}]\label{lem:integral}
Let $p=2$ or $p\equiv3\pmod4$. If $x,y\in\Q_p$ and $x^2+y^2=1$, then $x,y\in\Z_p$, and $(x,y)\bmod p$ is a
nonzero point of the circle $C_p=\{(a,b)\in\F_p^2: a^2+b^2=1\}$.
\end{lemma}

\begin{proof}
Suppose $-k=\min(v_p(x),v_p(y))<0$. If only one of $x,y$ had valuation $-k$, then $v_p(x^2+y^2)=-2k\ne0$. So
$x=p^{-k}a$, $y=p^{-k}b$ with $a,b\in\Z_p^\times$ and $a^2+b^2=p^{2k}$. For odd $p$ this gives
$(a/b)^2\equiv-1\pmod p$, impossible for $p\equiv3\pmod4$. For $p=2$, odd squares are $1\bmod8$, so
$v_2(a^2+b^2)=1$, while $v_2(p^{2k})$ is even. Hence $x,y\in\Z_p$, and reducing $x^2+y^2=1$ gives the rest.
\end{proof}

\begin{proposition}[cf.\ {\cite[Remark~3.6]{Madore}}]\label{prop:upper}
If $p=2$ or $p\equiv3\pmod4$, then $\chiB(\Q_p^2)\le\chi(\F_p^2)$, by a locally constant colouring.
\end{proposition}

\begin{proof}
Let $r(v)$ be the representative of $v+\Z_p^2$ whose coordinates are sums $\sum_{j<0}a_jp^j$ with
$0\le a_j<p$; it is locally constant. Given a proper colouring $\kappa$ of $\F_p^2$, put
$c(v)=\kappa\bigl((v-r(v))\bmod p\bigr)$. If $v$ and $w$ are adjacent, $v-w\in\Z_p^2$ by
Lemma~\ref{lem:integral}, so $r(v)=r(w)$, and $(v-r(v))-(w-r(w))=v-w$ reduces to a point of $C_p$; hence
$c(v)\ne c(w)$. The colouring is locally constant, so it is Borel.
\end{proof}

For $p=2$, $C_2=\{(1,0),(0,1)\}$ and $\kappa(a,b)=a+b$ is a proper $2$-colouring; for $p=3$,
$C_3=\{(\pm1,0),(0,\pm1)\}$ and $\kappa(a,b)=a+b\bmod3$ is a proper $3$-colouring. For $p=7$ the graph
$\F_7^2$ is $4$-colourable (Moorhouse \cite{Moorhouse}, Vinh \cite{Vinh}; a colouring is stored in \cite{Repo}).

\section{Measurable colourings: proof of \texorpdfstring{Theorem~\ref{thm:A}}{Theorem 1}}

The following lemma is the ratio bound of Bachoc, DeCorte, Oliveira and Vallentin \cite{BDOV}, which
\cite[Theorem~2.5]{BMK} states on $\Q_p^n$; we include the short proof for completeness. The new input is
Lemma~\ref{lem:levels}.

\begin{lemma}[Delsarte--Hoffman on a compact group]\label{lem:DH}
Let $\Gamma$ be a compact abelian group with Haar probability measure $\mu$, and $\sigma$ a probability measure
on a closed set $S\subset\Gamma$ with $S=-S$, $0\notin S$ and $\sigma$ invariant under $s\mapsto-s$. For a
character $\psi$ of $\Gamma$ put $\hat\sigma(\psi)=\int\psi\,d\sigma$, which is real, and let
$m=\inf\{\hat\sigma(\psi):\psi\ne1\}$, the infimum over the nontrivial characters; suppose $m<0$. If
$A\subset\Gamma$ is measurable and $A\cap(A+s)=\emptyset$ for every $s\in S$, then $\mu(A)\le -m/(1-m)$.
\end{lemma}

\begin{proof}
The function $f(x)=\mu(A\cap(A-x))$ is continuous and positive definite, with $\hat f(\psi)=|\hat 1_A(\psi)|^2\ge0$,
$\hat f(1)=\mu(A)^2$ and $\sum_\psi\hat f(\psi)=\|1_A\|_2^2=\mu(A)$ (Parseval); so $f=\sum_\psi\hat f(\psi)\psi$
uniformly. By hypothesis $f=0$ on $S$, hence
\[
0=\int f\,d\sigma=\sum_\psi\hat f(\psi)\hat\sigma(\psi)\ge\mu(A)^2+m\bigl(\mu(A)-\mu(A)^2\bigr),
\]
which gives the bound.
\end{proof}

Let $p\equiv3\pmod4$. Identify $\Z_p^2$ with the ring $\Oo=\Z_p[i]$ of integers of the unramified quadratic
extension $E=\Q_p(i)$ by $(x,y)\mapsto x+iy$; then $Q$ becomes the norm $\N z=z\bar z$, and the unit circle is the
compact group $S=\{z\in\Oo:\N z=1\}$, with its Haar probability measure $\sigma$. Since $E/\Q_p$ is unramified and
$p$ is odd, the norm maps $\Oo^\times$ onto $\Z_p^\times$ and $1+p^r\Oo$ onto $1+p^r\Z_p$ for every $r\ge1$; so
the image of $S$ in $(\Oo/p^r)^\times$ is the whole norm-one subgroup, of order $(p+1)p^{r-1}$, and the classes of
$S$ modulo $p^r$ all have measure $1/((p+1)p^{r-1})$. Write $e(t)=\exp(2\pi it)$ and $e_p(t)=e(\{t\}_p)$ for
$t\in\Q_p$. The characters of $\Oo$ are $\psi_\xi(z)=e_p(\Tr(\xi z))$ for $\xi\in E/\Oo$ (the different of $E/\Q_p$
is trivial), the trivial one being $\xi=0$; call $\ell=-v_p(\xi)$ the \emph{level} of $\psi_\xi$. Since $-1\in S$
and $\bar S=S$, every $\hat\sigma(\xi)=\int_S\psi_\xi\,d\sigma$ is real.

\begin{lemma}\label{lem:levels}
Let $\xi=p^{-\ell}\xi_0$ with $\xi_0\in\Oo^\times$.
\begin{enumerate}
\item If $\ell=1$, then $\hat\sigma(\xi)=\lambda/(p+1)$, where $\lambda$ is an eigenvalue of the graph
$\F_p^2$; more precisely $\lambda=-K(n)$ with $n=\N\xi_0\bmod p$ and the Kloosterman sum
$K(n)=\sum_{x\in\F_p^\times}e\bigl((x+n/x)/p\bigr)$, so $|\lambda|\le2\sqrt p$.
\item If $\ell\ge2$, then $|\hat\sigma(\xi)|\le 2/\bigl((p+1)p^{\lfloor\ell/2\rfloor-1}\bigr)\le2/(p+1)$.
\end{enumerate}
\end{lemma}

\begin{proof}
(1) The character factors through $\Oo/p\Oo=\F_{p^2}$, and $S$ maps onto the norm-one group
$\bar C=\{z\in\F_{p^2}:z\bar z=1\}$, of order $p+1$, with fibres of equal measure. So
$\hat\sigma(\xi)=\frac1{p+1}\sum_{z\in\bar C}\psi(\Tr(cz))$ with $c=\xi_0\bmod p$ and $\psi(t)=e(t/p)$, the
eigenvalue of $\F_p^2$ at the corresponding character. That these eigenvalues are Kloosterman sums is due to
Medrano, Myers, Stark and Terras \cite{MMST}; here is the short argument. With $n=\N c\ne0$, $w=cz$ gives
$\sum_{z\in\bar C}\psi(\Tr(cz))=\sum_{\N w=n}\psi(\Tr w)$. An element $w$ with $\N w=n$ and $\Tr w=t$ is a root of
$X^2-tX+n$, and conversely each root outside $\F_p$ has this norm and trace, while a root $w\in\F_p$ has norm
$w^2$ and trace $2w$, which fits only when $t^2=4n$. So there are $1-\eta(t^2-4n)$ such $w$, where $\eta$ is the
quadratic character of $\F_p$. Similarly $x+n/x=t$ has $1+\eta(t^2-4n)$ solutions $x\in\F_p^\times$. As
$\sum_t\psi(t)=0$, the first sum is $-\sum_t\psi(t)\eta(t^2-4n)=-K(n)$, and Weil's bound \cite{Weil} gives
$|K(n)|\le2\sqrt p$.

(2) Let $j=\lceil\ell/2\rceil$, $r=\ell-j=\lfloor\ell/2\rfloor\ge1$, and $S_j=\{s\in S:s\equiv1\bmod p^j\}$. For
$s_0\in S$, the function $s\mapsto\psi_\xi(s_0s)/\psi_\xi(s_0)=e_p(\Tr(\xi s_0(s-1)))$ is a character of $S_j$,
because $(1+p^jt)(1+p^jt')=1+p^j(t+t')+p^{2j}tt'$ and $v_p(\xi p^{2j})\ge0$. Its average over $S_j$ is $0$
unless it is trivial. The elements $(1+p^jia/2)/(1-p^jia/2)$, $a\in\Z_p$, lie in $S_j$ and are
$\equiv1+p^jia\pmod{p^{2j}}$, so triviality forces $e_p(\Tr(i\,p^{j-\ell}a\,\xi_0s_0))=1$ for all $a\in\Z_p$, that
is, $\Tr(i\xi_0s_0)=-2\operatorname{Im}(\xi_0s_0)\equiv0\pmod{p^r}$. Then $\xi_0s_0\equiv u\pmod{p^r}$ for some
$u\in\Z_p$ with $u^2\equiv\N\xi_0\pmod{p^r}$, so $s_0$ lies in one of at most two classes modulo $p^r$, each of
measure $1/((p+1)p^{r-1})$; and the cosets of $S_j$ refine these classes since $j\ge r$. Averaging over the cosets
of $S_j$ first gives the bound.
\end{proof}

\begin{proof}[Proof of Theorem~\ref{thm:A}]
The upper bounds are Proposition~\ref{prop:upper} and \cite[Theorem~1]{Vinh}. For the lower bound, let $c$ be a
proper colouring of $\Q_p^2$ with $k$ measurable classes. Their traces on $\Z_p^2$ cover it, so one of them, $A$,
has $\mu(A)\ge1/k$, and since $S\subset\Z_p^2$ no two points of $A$ differ by an element of $S$. By
Lemma~\ref{lem:levels}, $m\ge-\max(2\sqrt p,2)/(p+1)=-2\sqrt p/(p+1)$, and $m<0$ because the eigenvalues of
$\F_p^2$ sum to $0$; so Lemma~\ref{lem:DH} gives $1/k\le\mu(A)\le|m|/(1+|m|)$, that is,
$k\ge1+1/|m|\ge1+(p+1)/(2\sqrt p)$.
\end{proof}

With the exact least eigenvalue $\lambda_{\min}$ of $\F_p^2$, the same proof gives
$\chiM(\Q_p^2)\ge1+(p+1)/\max(|\lambda_{\min}|,2)$. We computed $\lambda_{\min}$ in interval arithmetic
(\texttt{padic\_measurable.py} in \cite{Repo}); Table~\ref{tab:meas} shows the result. Weil's bound alone gives
at least $6$ colours for $p\ge62$, $7$ for $p\ge98$, $8$ for $p\ge142$ and $9$ for $p\ge194$; with the exact
spectra (computed for $p<200$), $\chiM(\Q_p^2)\ge6$ already for $p=59$, $\ge7$ for $p=71$, $\ge8$ for every
$p\ge131$, and $\ge9$ for $p=191$.

\begin{table}[ht]
\centering
\begin{tabular}{rrrr}
$p$ & $\lambda_{\min}$ & $1+1/|m|$ & $\chiM(\Q_p^2)\ge$\\\hline
3 & $-2.0000$ & $3.0000$ & 3\\
7 & $-4.4940$ & $2.7802$ & 3\\
11 & $-4.7958$ & $3.5022$ & 4\\
19 & $-7.6097$ & $3.6282$ & 4\\
23 & $-7.9607$ & $4.0148$ & 5\\
31 & $-10.6665$ & $4.000044$ & 5\\
43 & $-11.9622$ & $4.6782$ & 5\\
47 & $-12.8465$ & $4.7364$ & 5\\
59 & $-14.9176$ & $5.0221$ & 6\\
67 & $-15.0567$ & $5.5163$ & 6\\
71 & $-14.1728$ & $6.0801$ & 7\\
79 & $-17.1019$ & $5.6778$ & 6\\
83 & $-17.7087$ & $5.7434$ & 6\\
103 & $-20.0322$ & $6.1916$ & 7\\
107 & $-19.0495$ & $6.6694$ & 7\\
\end{tabular}
\caption{Lower bounds for the measurable chromatic number of $\Q_p^2$, $p\equiv3\pmod4$, from the least eigenvalue
of $\F_p^2$ (rounded; the bounds are rigorous lower ends of 120-bit intervals).}
\label{tab:meas}
\end{table}

Together with Theorem~\ref{thm:B}, $\chiM(\Q_p^2)=\chiB(\Q_p^2)=\chi(\Q_p^2)$ for $p=2,3,7$, while
$4\le\chiM(\Q_p^2)\le\chiB(\Q_p^2)\le5$ for $p=11,19$, where $\chi(\F_p^2)=5$ (\cite{Repo} stores $5$-colourings of
$\F_{11}^2$ and $\F_{19}^2$, and certificates that $4$ colours do not suffice).

\section{Lower bounds from number fields}

If $K$ is a field and $\iota:K\hookrightarrow\Q_p$, then $(x,y)\mapsto(\iota x,\iota y)$ is an injective
homomorphism $G(K)\to G(\Q_p)$, so $\chi(\Q_p^2)\ge\chi(K^2)$. For odd $p$ and $p\nmid d$, Hensel's lemma gives
$\Q(\sqrt d)\subset\Q_p$ exactly when $d$ is a square modulo $p$.

\begin{proof}[Proof of Theorem~\ref{thm:B}]
The upper bounds are Proposition~\ref{prop:upper}. The graph has edges, so $\chi(\Q_2^2)\ge2$. For $p=3$, besides
Madore's $9$-cycle, the unit vectors
\[
(1,0),\quad(0,1),\quad\Bigl(\tfrac{-1+\sqrt7}4,\tfrac{-1-\sqrt7}4\Bigr),\quad
\Bigl(-\tfrac{\sqrt7}4,-\tfrac34\Bigr),\quad\Bigl(-\tfrac34,\tfrac{\sqrt7}4\Bigr)
\]
sum to $0$, and their partial sums $0$, $(1,0)$, $(1,1)$, $(\frac{3+\sqrt7}4,\frac{3-\sqrt7}4)$,
$(\frac34,-\frac{\sqrt7}4)$ are distinct, so $G(\Q(\sqrt7))$ contains a $5$-cycle; as $7\equiv1\pmod3$,
$\sqrt7\in\Q_3$ and $\chi(\Q_3^2)\ge3$. Finally $11\equiv2^2\pmod7$, so $\sqrt{11}\in\Q_7$, and the
$76$-vertex unit-distance graph over $\Q(\sqrt{11})$ of \cite{QP} has no proper $3$-colouring (a DRAT proof
checked by \texttt{drat-trim} and by the verified checker \texttt{cake\_lpr}, and a proof in Lean); so
$\chi(\Q_7^2)\ge4$.
\end{proof}

\begin{proof}[Proof of Theorem~\ref{thm:D}]
Let
\begin{gather*}
\mathcal D=\{11,23,35,47,59,71,95,119,131,155,179,191,239,251,\\
263,359,431,443,455,491,599,611,791,851,911,935,959\}.
\end{gather*}
For each $d\in\mathcal D$, \cite{QP} gives a unit-distance graph over $\Q(\sqrt d)$ with no proper
$3$-colouring. A segmented sieve with Jacobi symbols (\texttt{padic\_reach.c} in \cite{Repo}, about half a
minute) shows that for every prime $p\equiv3\pmod4$ with $7\le p<2\,129\,503\,819$ some $d\in\mathcal D$ is a
nonzero square modulo $p$, so $\Q(\sqrt d)\subset\Q_p$ and $\chi(\Q_p^2)\ge4$; a separate script checks every
prime below $10^7$ again with Euler's criterion. The prime $2\,129\,503\,819$ is the first after $3$ at which every
$d\in\mathcal D$ is a non-square. For $p=83$, the numbers $3$, $11$ and $247$ are squares modulo $83$, and
\cite{Repo} contains a $5$-chromatic unit-distance graph over $\Q(\sqrt3,\sqrt{11},\sqrt{247})$.
\end{proof}

\begin{proof}[Proof of Proposition~\ref{prop:E}]
Let $M$ be the Galois closure of the compositum of $\Q(i)$ and the fields in $\mathcal K$, $G=\operatorname{Gal}(M/\Q)$,
$D\subset G$ the decomposition group of a prime of $M$ above $3$, $I\subset D$ its inertia group, and $\sigma\in D$
an element mapping to the Frobenius generator of $D/I$, so that $D=\langle\sigma\rangle I$. By Chebotarev's
theorem infinitely many primes $p$, unramified in $M$, have Frobenius conjugate to $\sigma$. Since $3$ is inert in
$\Q(i)$, $\sigma$ acts on the residue field $\F_9$ of $\Q(i)$ at $3$ by $x\mapsto x^3$ and so moves $i$; hence
$p\equiv3\pmod4$. Suppose $K=M^H\in\mathcal K$ embeds in $\Q_p$. Then a prime of $K$ above $p$ has residue degree
$1$, so the Frobenius at $p$, and with it $\sigma$, which has the same cycle type on $G/H$, fixes a coset $gH$;
and $D\cdot gH=I\langle\sigma\rangle\cdot gH=I\cdot gH$. The primes of $K$ above $3$ correspond to the $D$-orbits
on $G/H$, with residue degree the number of $I$-orbits in the orbit (see \cite[Chapter~I, \S9]{Neukirch}); so the
prime given by $D\cdot gH$ has residue field $\F_3$, in which $-1$ is not a square. The argument of
Lemma~\ref{lem:integral} and Proposition~\ref{prop:upper}, with a uniformiser in place of $p$ and
$\kappa(a,b)=a+b\bmod 3$, then gives $\chi(K^2)\le3$, a contradiction.
\end{proof}

For quadratic fields this is Fischer's theorem that $\chi(\Q(\sqrt N)^2)\le3$ when $N\equiv0,1\pmod3$
\cite{Fischer1990}. A uniform bound $\chi(\Q_p^2)\ge4$ for all $p\ge7$ would follow, for example, from
$\chi(\Q(\sqrt q)^2)\ge4$ for every prime $q\equiv11\pmod{12}$: by Dirichlet's theorem every prime
$p\equiv3\pmod4$ has such a $q$ with $q\equiv1\pmod p$, so that $\Q(\sqrt q)\subset\Q_p$.

\section{The isotropic case}

\begin{proof}[Proof of Proposition~\ref{prop:C}]
Let $HG(F)$ be the graph on $F^2$ in which $(x,y)$ and $(x',y')$ are adjacent when $(x-x')(y-y')=1$. By
\cite[Lemma~2.7]{BMK}, $HG(F)$ is an induced subgraph of the graph of any isotropic form over $F$, and it contains
$HG(\Q)$. The map $(x,y)\mapsto(x-y,x+y)$ sends the graph of $x^2-y^2$ on $\Q^2$ onto $HG(\Q)$, because
$(\Delta x)^2-(\Delta y)^2=(\Delta x-\Delta y)(\Delta x+\Delta y)$. Davies \cite[Theorem~3]{Davies}, from a
theorem of Graham, proved that every finite colouring of $\Q^2$ has a monochromatic pair with
$(x-x')^2-(y-y')^2=1$. Hence no finite colouring is proper. The form $x^2+y^2$ is isotropic over $F$ exactly when
$-1$ is a square in $F$, and over $\Q_p$ this happens exactly when $p\equiv1\pmod4$.
\end{proof}

Madore had observed that $\chi(\R^{1,1})=\infty$ would give $\chi(\C^2)=\infty$ \cite[\P6.4]{Madore}, and Davies
recorded this consequence of his theorem; the proposition is the same remark for any field of characteristic
$0$.

\section{Questions}

\begin{enumerate}
\item For $p\equiv3\pmod4$, $\chiB(\Q_p^2)\to\infty$, but the ordinary $\chi(\Q_p^2)$ is only known to be at
least $4$ for $7\le p<2\,129\,503\,819$ (Theorem~\ref{thm:D}) and at least $5$ at $p=83$. Is it bounded? A finite
subgraph whose coordinates generate a field with a real embedding lies in $\R^2$, so more than seven colours would
need a field with no real embedding. If $\chi(K^2)$ is unbounded over number fields $K$ in which $-1$ is not a
square, then so is $\chi(\Q_p^2)$ over $p\equiv3\pmod4$, since by Chebotarev's theorem each such $K$ embeds in
$\Q_p$ for infinitely many $p\equiv3\pmod4$. If $\chi(\Q_p^2)$ is bounded, the $p$-adic planes are a quantitative
analogue, along a family, of the gap between ordinary and Borel chromatic numbers that \cite[Remark~1.3]{BMK}
asks about (whose own instance, for indefinite real forms, is ruled out by Davies's theorem).
\item Is $\chi(\Q_{11}^2)$ equal to $4$ or $5$? Since $47$, $3$ and $5$ are squares modulo $11$, the field
$\Q_{11}$ contains $\Q(\sqrt{47})$ and $\Q(\sqrt3,\sqrt5)$; so $\chi(\Q_{11}^2)=4$ would give
$\chi(\Q(\sqrt{47})^2)=4$, a case Moorhouse asked about \cite{MoorhouseTalk} that is still open \cite{QP}, and
$\chi(\Q(\sqrt3,\sqrt5)^2)\le4$. The colourings of $\Q_{11}^2$ that are constant on the cosets of $11^k\Z_{11}^2$
need five colours for $k=1,2,3$ \cite{Repo}. At $p=19$ the question about $\Q(\sqrt{47})$ is the same, as
$47\equiv3^2\pmod{19}$, but $\Q(\sqrt3,\sqrt5)$ does not embed in $\Q_{19}$; there we only know that the
colourings constant on the cosets of $19\Z_{19}^2$ need five colours.
\item How fast does $\chiB(\Q_p^2)$ grow? Theorem~\ref{thm:A} places it between $(\frac12+o(1))\sqrt p$ and
$(p+1)/2$.
\end{enumerate}

\begin{thebibliography}{99}
\bibitem{BDOV} C. Bachoc, E. DeCorte, F. M. de Oliveira Filho and F. Vallentin, \emph{Spectral bounds for the
independence ratio and the chromatic number of an operator}, Israel J. Math. \textbf{202} (2014), 227--254.
\bibitem{BMK} M. Bardestani and K. Mallahi-Karai, \emph{On a generalization of the Hadwiger--Nelson problem},
Israel J. Math. \textbf{217} (2017), 313--335; arXiv:1507.05300 (we quote the arXiv version).
\bibitem{BMK2} M. Bardestani and K. Mallahi-Karai, \emph{On the chromatic number of structured Cayley graphs},
arXiv:1511.02427 (2015, revised 2018).
\bibitem{BMK3} M. Bardestani and K. Mallahi-Karai, \emph{Polynomial configurations in sets of positive upper
density over local fields}, J. Anal. Math. (2020); arXiv:1701.06024.
\bibitem{Davies} J. Davies, \emph{Chromatic number of spacetime}, arXiv:2308.16885 (2023, revised 2024).
\bibitem{deGrey} A. D. N. J. de Grey, \emph{The chromatic number of the plane is at least 5}, Geombinatorics
\textbf{28} (2018), 18--31; arXiv:1804.02385.
\bibitem{EI} G. Exoo and D. Ismailescu, \emph{The chromatic number of the plane is at least $5$: a new proof},
Discrete Comput. Geom. \textbf{64} (2020), 216--226.
\bibitem{Fischer1990} K. G. Fischer, \emph{Additive $K$-colorable extensions of the rational plane}, Discrete
Math. \textbf{82} (1990), 181--195.
\bibitem{Repo} S. Gal\'an, \emph{The Hadwiger--Nelson problem over number fields}, code, data and formal
proofs, 2026. Archived at Zenodo, \href{https://doi.org/10.5281/zenodo.22976635}{doi:10.5281/zenodo.22976635}
(all versions); source at \url{https://github.com/decalion89/chromatic-number-of-the-plane}; the scripts cited here
are in \texttt{data/quadratic\_planes/scripts/}, and the account is \texttt{notes/quadratic\_planes.md}, \S5.
\bibitem{QP} S. Gal\'an, \emph{Real quadratic planes that need four colours}, draft, 2026, in \cite{Repo}
(\texttt{papers/quadratic-planes/}).
\bibitem{Johnson1987} P. D. Johnson Jr., \emph{Two-colorings of real quadratic extensions of $\Q^2$ that forbid
many distances}, Congr. Numer. \textbf{60} (1987), 51--58.
\bibitem{Madore} D. A. Madore, \emph{The Hadwiger--Nelson problem over certain fields}, arXiv:1509.07023 (2015).
\bibitem{mathlib} The mathlib Community, \emph{The Lean mathematical library}, CPP 2020, 367--381;
arXiv:1910.09336.
\bibitem{Lean} L. de Moura and S. Ullrich, \emph{The Lean 4 theorem prover and programming language}, CADE 28,
Lecture Notes in Comput. Sci. \textbf{12699} (2021), 625--635.
\bibitem{MMST} A. Medrano, P. Myers, H. M. Stark and A. Terras, \emph{Finite analogues of Euclidean space},
J. Comput. Appl. Math. \textbf{68} (1996), 221--238.
\bibitem{Moorhouse} G. E. Moorhouse, \emph{On the chromatic numbers of planes}, preliminary draft, revised
3 March 2010, \url{https://www.ericmoorhouse.org/pub/chromatic.pdf}.
\bibitem{MoorhouseTalk} G. E. Moorhouse, \emph{Colouring the plane}, talk at Designs, Codes \& Geometries
2010, slides at \url{https://www.ericmoorhouse.org/slides/ebertfest.pdf}.
\bibitem{Neukirch} J. Neukirch, \emph{Algebraic number theory}, Grundlehren Math. Wiss. \textbf{322}, Springer,
Berlin, 1999.
\bibitem{Soifer} A. Soifer, \emph{The mathematical coloring book}, Springer, New York, 2009.
\bibitem{Vinh} L. A. Vinh, \emph{On chromatic number of unit-quadrance graphs (finite Euclidean graphs)},
arXiv:math/0510092 (2005).
\bibitem{Weil} A. Weil, \emph{On some exponential sums}, Proc. Nat. Acad. Sci. U.S.A. \textbf{34} (1948),
204--207.
\bibitem{Woodall} D. R. Woodall, \emph{Distances realized by sets covering the plane}, J. Combin. Theory Ser. A
\textbf{14} (1973), 187--200.
\end{thebibliography}

\end{document}
